In performing this study it became pertinent to understand where
  the sources of the muons under consideration were originating.

A large set of Pythia events were generated to understand the sources, namely,
  101 million raw events of \textit{HardQCD:hardbbbar}, 
  576 million raw events for \textit{HardQCD:hardccbar},
  and 576 million raw events for \textit{HardQCD:all}.
Imposing the requirement of at least two muons with 
  a $\pt$ greater than 4.0~\GeV within $|\eta|<2.5$,
  on these Pythia events, resulted in 59k, 9.1k, and 4.0k 
  dimuon pairs respectively.

The overwhelming brunt of the pairs were generated via the 
  hard scattering process and not from the initial parton shower.
Pairs were created both via pair production and flavor excitation.
The origins of these pairs can be seen more cleanly in Figure~\ref{fig:Parents}.

\begin{figure}[h]
\centering
\includegraphics[width=0.5\textwidth]{figures/MuonSources/parents}
\caption{
   Parents vs governing process and origin (hard scattering vs initial parton shower) of Pythia-generated dimuon pairs. 
   b: b-quark, c: c-quark, g: gluon, q: other quark, hf: heavy flavor, onia: quarkonia, ex: excitation, hs: hard scattering,
   o11: both muons coming from hard scattering,
   o12: one muon from hard scattering and one muon from the initial parton shower,
   o22: both muons from the initial parton shower.
   For examples "c ex" means c-flavor excitation process.
   "gg/qq$\rightarrow$bb" would signify a b pair production process,
   whereas a Parent "bbc" would indicate one muon from a b-decay, the other from a b$\rightarrow$c decay.
\label{fig:Parents}
} 
\end{figure}

\begin{table}[h]
	\centering
	\begin{tabularx}{0.65\textwidth}{ | X | X | X |  }
	\hline
		HardQCD Process & $c\bar{c}$ $\sigma$ [$\mu b$]& $b\bar{b}$ $\sigma$ [$\mu b$] \\ \hline
		$c\bar{c}$ & 150 $\pm$ 0.5 & N/A \\ \hline
		$b\bar{b}$ & N/A & 46 $\pm$ 0.2  \\ \hline 
 		all & 167 $\pm$ 10 & 49 $\pm$ 5 \\	\hline
	\end{tabularx}
	\caption{Cross-sections of $c\bar{c}$, $b\bar{b}$ for HardQCD:hardccbar, HardQCD:hardbbbar, and HardQCD:all with no dimuon quantity or kinematic restrictions}
	\label{tab:AllEvents}
\end{table}

      
We also ran these simulations with no dimuon requirements, and we see in Table ~\ref{tab:AllEvents} that, as we might expect,
  the cross-section for muons created is about three times greater for $c\bar{c}$ than $b\bar{b}$ when including all events
  with no dimuon restriction or kinematic cuts.
However, perhaps somewhat surprisingly, although hinted at by the number of pairs generated per events ran, and Figure~\ref{fig:Parents} being dominated 
  by bb parents coming from both b-flavor excitation and pair production, we see more clearly in Figure ~\ref{fig:Pairs_bbvscc}, 
  the $\Delta\phi$ distributions are plotted for normalized yields of both $c\bar{c}$ and $b\bar{b}$ events, that given the requirements (kinematic,
  and more notably requiring two muons) the yields in the away-side peak of interest from $b\bar{b}$ events are close to
  10 times the yields from the $c\bar{c}$ events. This is also reflected in the cross-section shown in Table ~\ref{tab:DimuonsReq} when the dimuon requirements are employed.


\begin{figure}[h]
\centering
\includegraphics[width=0.5\textwidth]{figures/MuonSources/pairs_bbvscc}
\caption{
   Scaled yields of all dimuon pairs from all possible decays channels for both 
   $c\bar{c}$ (red) and $b\bar{b}$ (blue) with the kinematic requirements enforced.	
   No \Deta\ cuts are imposed on the muons.
} \label{fig:Pairs_bbvscc}
\end{figure}

\begin{table}[h]
\centering
\begin{tabularx}{0.65\textwidth}{ | X | X | X |  }
\hline
	HardQCD Process & $c\bar{c}$  $\sigma$ [$\mu b$]& $b\bar{b}$ $\sigma$ [$\mu b$]\\ \hline
	$c\bar{c}$ & 0.0021 $\pm$ 2e-5 & N/A \\ \hline
	$b\bar{b}$ & N/A & 0.018 $\pm$ 1e-4  \\ \hline 
		all & 0.0023 $\pm$ 3e-5 & 0.020 $\pm$ 2e-4 \\	\hline
\end{tabularx}
\caption{
  Cross-sections of $c\bar{c}$, $b\bar{b}$ for HardQCD:hardccbar, 
  HardQCD:hardbbbar, and HardQCD:all with listed dimuon quantity or kinematic restrictions
 \label{tab:DimuonsReq}
}
\end{table}


Additionally, these simulated muon pairs were split into same sign and opposite sign pairs.
The scaled yields are shown in Figure ~\ref{fig:ss_v_os_yields}. We can see that opposite-sign
  pairs should out-weigh same-sign pairs about 2:1 in the away-side peak for $bb$ pairs, 
  and $c\bar{c}$ pairs can only exist in the opposite-sign configuration, as expected.

\begin{figure}[h]
\centering
\includegraphics[width=0.9\textwidth]{figures/MuonSources/ss_v_os_yields}
\caption{
   Same-sign and Opposite-sign scaled yields of all dimuon pairs from all possible decays channels for both 
   $c\bar{c}$ (red) and $b\bar{b}$ (blue) with the kinematic requirements enforced.	
   No \Deta\ cuts are imposed on the muons.
} \label{fig:ss_v_os_yields}
\end{figure}


\begin{comment}
I don't think this is needed anymore and can be deleted.

Additionally, in for Pb+Pb and pp data, the yields were split into same sign pairs and opposite sign pairs.
FIGURE plots the Pb+Pb in points as a function of centrality and pp as the dotted line for same-sign pairs
  in the left most plot, opposite sign pairs in the central most plot and their ratio in the right plot.
While opposite sign pairs can be generated from from either $b\bar{b}$ or $c\bar{c}$ decays,
  but same sign pairs can only come from $b\bar{b}$ pairs via a $b\rightarrow c \rightarrow \mu$ decay.
We see here that about a third of the events are same-sign pairs, which come exclusively from bbbar events,
  and we expect at least $50\%$ of the opposite-sign pairs to come from bbbar events, although
  likely more given the Pythia findings.
Regardless however, we once again find that dimuon pairs with these requirements should be dominated by
  processes originating from bbbar events.
\end{comment}
